% !TeX root = ../main.tex

\section{Considerações finais}

A análise da vegetação nativa declarada no Cadastro Ambiental Rural (CAR) na bacia hidrográfica do rio Gramame permitiu identificar 8.224,30 ha de vegetação declarada espacialmente distinta, distribuídos em 279 polígonos associados ao recorte de 1.974 imóveis rurais. A identificação de sobreposições entre os polígonos demonstrou a importância do tratamento geométrico para evitar a dupla contagem das áreas.

A comparação com o MapBiomas Coleção 11, referente ao ano de 2025, apresentou índice de convergência espacial de 63,66\%. A formação florestal predominou entre as classes naturais coincidentes, enquanto pastagem e mosaico de usos concentraram a maior parte das divergências. Também foram identificados 3.012,30 ha de cobertura natural dentro dos imóveis rurais, mas fora dos polígonos de vegetação nativa declarada analisados.

Os resultados demonstram que a integração entre o CAR e o MapBiomas pode contribuir para a identificação de divergências espaciais e apoiar o direcionamento de análises ambientais mais detalhadas. Entretanto, essas diferenças não devem ser interpretadas isoladamente como irregularidades cadastrais ou ambientais, considerando as limitações relacionadas às datas de referência, à resolução espacial e aos procedimentos de classificação das bases utilizadas.

Para trabalhos futuros, recomenda-se incorporar análises temporais da cobertura da terra, imagens de maior resolução e avaliações de sensibilidade dos critérios de comparação, buscando compreender melhor as divergências identificadas.


\label{pag:fim-texto}
\newpage
